\section{Technical Research Lines}
\label{sec:methods}
A useful taxonomy has at least two independent layers: the graphic-design problem being solved and the computational mechanism used. The same representation may be produced by diffusion, autoregressive models, optimization, or multi-stage agents, while the same mechanism may produce several kinds of artifact.

\subsection{Constraint-Aware Layout Generation}
Layout generation predicts discrete elements and their geometric attributes. This includes unconstrained synthesis, category-conditioned generation, partial completion, and constraint-aware refinement. CGL-GAN conditions advertising layouts on visual context \citep{cglgan2022}; LayoutDETR applies detection-transformer-style conditioning to foreground composition \citep{layoutdetr2024}. LayoutDM and LayoutDiffusion illustrate diffusion over discrete layout attributes \citep{layoutdm2023,layoutdiffusion2023}, while LayoutFlow addresses continuous geometry \citep{layoutflow2024}. In comparing these systems, the allowed conditions and datasets must be matched: unconditional realism, completion, and relational constraint satisfaction are not equivalent tasks.

\subsection{Multi-Stage Graphic Design and Agents}
Integrated poster generation must coordinate language, asset production, layout decisions, and visual verification. COLE and OpenCOLE use hierarchical design decomposition \citep{cole2023,opencole2024}. CreatiDesign develops multi-conditional visual generation across primary content and decorative or textual components \citep{creatidesign2026}. Vinci combines data-driven poster creation with interactive editing feedback and user evaluations \citep{vinci2021}. Code-oriented systems increasingly produce the design document itself, as in DesignAsCode and Editable Visual Design \citep{designascode2026,editablevisual2026}; PSDesigner targets layered editor workflows \citep{psdesigner2026}. PosterAgent studies staged draft generation and refinement through agentic reinforcement learning \citep{posteragent2026}. We separate \emph{reflection during generation} from an edit request received after a completed artifact is delivered.

\subsection{Typography and Vector Synthesis}
Text rendering poses constraints not captured by photographic image fidelity. TextDiffuser makes text layout a first-class part of image generation \citep{textdiffuser2023}; GlyphDraw2 targets stylized, multilingual glyph-aware poster generation \citep{glyphdraw2025}. Word-As-Image optimizes letter shapes to communicate meaning while retaining legibility and font resemblance \citep{wordasimage2023}. De-Rendering Stylized Texts instead infers editable rendering parameters from a flattened text image \citep{derendering2021}. A replacement of semantic text, a modification of glyph outlines, and a reconstruction of font/effect parameters are three different editing operations.

Vector-centric work generates or embeds shape/path sequences. DeepSVG and SVGformer learn vector representations \citep{deepsvg2020,svgformer2023}, whereas VectorFusion and IconShop synthesize vectors conditioned on textual prompts \citep{vectorfusion2023,iconshop2023}. Their visual or semantic generation scores should not be interpreted as evidence for the broader class of targeted design edits.

\subsection{Reconstruction and Explicit Editing}
Another line starts from an existing artifact. LayerD decomposes rendered designs into layers \citep{layerd2025}; PosterCopilot reasons about composition and repeated layer-sensitive changes \citep{postercopilot2026}. MiLDEdit explicitly addresses instruction-driven modifications of multi-layer design documents, presenting a dedicated benchmark and evaluation framework \citep{mildedit2026}. PosterOmni spans local raster-image editing and global poster creation \citep{posteromni2026}. These tasks are meaningfully editing-oriented, but a pixel-space editing benchmark does not establish that the result is an editable PSD or DOM structure. Each study must be coded against its actual inputs, outputs, and scoring procedure.

\subsection{Beyond Poster-Specific Settings}
The same principles apply to scientific figures, graphical abstracts, and slide generation, although constraints such as factual correctness, visual rhetoric, and export format may differ from commercial advertisement design. A representation-centered classification permits such work to be compared without assuming that the metrics and data distributions are identical. Cross-domain synthesis should respect these differences rather than collapse them into a generic text-to-image category.
